\chapter{Integration into RTA-CAR Project}
\label{ch:IntGuide}

This chapter aims to provide a guideline for the integration of RTA-SUM SrvcMgr module into an RTA-CAR project.

Once RTA-SUM plugin is installed, RTA-SUM SrvcMgr gets available for being integrated into a project. Create a new RTA-CAR project in ISOLAR-AB and select the SDV 2.0 version. 
Open the RTA Code Generator dialog and enable RTA-SUM SrvcMgr by choosing the SDV 2.0 selected during the project creation. If the option Not Configured is chosen instead, the module will not be integrated.

Refer to RTA-SUM plugin Getting Started Guide for further details.

\section {Prerequisites}
\begin{itemize}
\item It is assumed the integrator has followed the RTA-SUM Getting Started Guide to install the RTA-SUM plugin.
\item In order to successfully follow this guide, you shall have the RTA-CAR toolchain installed and you must be familiar with the AUTOSAR specifications, terminology and methodology. You also must be familiar with the GM specification and GM-SUM modules functionalities.
\end{itemize}

\section{Create SrvcMgr module}

From the 'ECU Navigator' view, expand the Bsw Modules, and add the 'SrvcMgr' module to the 'Other Modules', as shown below:

%\begin{figure}[H]
%	\centering
%	\includegraphics[scale=0.35]{images/createSSM.png}
%\end{figure}

\section{SrvcMgr Configuration}
Once RTA-SUM plugin has been installed and the RTA-CAR project has been created, the module SrvcMgr should
appear. From here you can open and begin configuring the module. Specific description of each configurable
parameter is available in the “Description” view of ISOLAR-AB when a parameter is selected.
Please refer to Section 3 for more details about RTA-SUM SrvcMgr configuration.
